\section{Related Work}
\label{sec:related_work}

\noindent{\bf Memory-based Robotic Manipulation Benchmarks.}
Most prior robotic manipulation benchmarks \cite{liu2023libero, li24simpler, mees2022calvin, zhu2020robosuite} implicitly involve temporal memory but rarely require it: policies conditioned only on the current observation can already achieve high success rates \cite{pi05, wang2025vlaadaptor, dai2025aimbot}. This suggests that success often relies on local perception rather than true history-based reasoning. For example, object states are continually updated during online execution, allowing the agent to infer the next subtask without relying on past observations.
MemoryBench \cite{fang2025sam2act} is the first attempt to introduce explicit memory requirements, focusing on spatial location recall, but its tasks remain simple and nearly solved. MIKASA-Robo \cite{cherepanov2025memory}, originally designed for reinforcement learning, includes some memory-related tasks, yet these tasks remain short in horizon, are limited in long-term dependencies, and lack sufficient high-quality demonstrations for effective imitation learning.
Recent memory-based policies \cite{shi2025memoryvla, sridhar2025memer} are typically evaluated on narrow, self-designed tasks, hindering systematic comparison and understanding.
In contrast, \benchname\xspace provides a challenging testbed with \emph{explicit} history dependence, with diverse memory types and sufficient demonstrations.

\noindent{\bf Memory-based Robotic Manipulation Models.}
Existing approaches can be broadly categorized according to their memory representations.
(1) \textit{Symbolic memory} relies on non-differentiable abstractions derived from interaction history or external modules. For example, HistRISE \cite{chen2025history} tracks object-centric 3D points as symbolic states,
UniVLA \cite{bu2025learning} incorporates temporal context by
adding past actions into input prompts.
More recent methods, such as MemER \cite{sridhar2025memer} and Gemini-Robotics-1.5 \cite{team2025geminirobot}, generate language-based subgoals with large vision-language models, offloading long-term memory reasoning to agentic pipelines, but introduce costly modular inference and preclude end-to-end optimization.
(2) \textit{Perceptual memory} represents history as differentiable visual or multimodal features. ContextVLA \cite{jang2025contextvla} appends past visual tokens directly to the transformer input as raw contextual tokens. In contrast, memory-bank approaches encode visual features together with auxiliary signals, such as task instructions \cite{shi2025memoryvla, koo2025hamlet} or action heatmaps \cite{fang2025sam2act}, and cache these multimodal embeddings for subsequent retrieval.
(3) \textit{Recurrent memory} compresses history into fixed-size latent states through iterative updates. Early work, such as BC-RNN \cite{mandlekar2021matters}, models temporal dependencies using recurrent neural networks. More recent approaches, including MITL \cite{zhou2025mtil} and RoboMamba \cite{liu2024robomamba}, adopt Mamba-style state-space models \cite{gu2024mamba} to better capture long-horizon dependencies.
%Despite showing the importance of memory, these approaches vary widely in architecture and evaluation, making systematic comparison difficult. 
%To address this gap, we construct a family of memory-augmented VLA models on the same backbone and evaluate diverse memory representations and integration strategies under a controlled, standardized setting.
To evaluate differnet memory designs, we construct a family of memory-augmented VLA models on the same backbone and evaluate diverse memory representations and integration strategies under a controlled, standardized setting.
